  \def\stat{kriv-lin}





\def\tit{СТАТИЧЕСКИЙ СПОСОБ СТЕГАНОГРАФИЧЕСКОГО 


ВСТРАИВАНИЯ ИНФОРМАЦИИ НА~ОСНОВЕ LSB}





\def\titkol{Статический способ стеганографического 


встраивания информации на~основе LSB}








\def\aut{И.\,А.~Кривошеев$^1$, М.\,А.~Линник$^2$}





\def\autkol{И.\,А.~Кривошеев, М.\,А.~Линник}





\titel{\tit}{\aut}{\autkol}{\titkol}





\index{Кривошеев И.\,А.}


\index{Линник М.\,А.}


\index{Krivosheev I.\,A.}


\index{Linnik M.\,A.}











%{\renewcommand{\thefootnote}{\fnsymbol{footnote}}


%\footnotetext[1]


%{Исследование выполнено за счет гранта Минобрнауки РФ по заявке №\,2020-1902-01-016.}} 








\renewcommand{\thefootnote}{\arabic{footnote}}


\footnotetext[1]{Институт горного дела Дальневосточного отделения Российской академии наук; Хабаровский Федеральный 


исследовательский центр Дальневосточного отделения Российской академии наук, \mbox{igork@as.khb.ru}}


\footnotetext[2]{Вычислительный центр Дальневосточного отделения Российской академии наук, 


\mbox{linnik.max1995@mail.ru}}





 \label{st\stat}





 \thispagestyle{headings}


 


 \vspace*{-12pt}


  


  


  


      


   


   \Abst{Предлагается новый подход к~сокрытию информации на основе метода LSB (least significant 


bit). Предложен статический способ встраивания информации в~цветное изображение, основанный на 


матричном представлении отдельных блоков и~оценке модуля их детерминанта. Численным 


моделированием была оценена возможность предложенного способа противостоять различным 


методам стегоанализа. Показано его преимущество как в~степени защищенности, так и~в объеме 


встраиваемой информации. Данный алгоритм может быть использован для встраивания информации в~изображения форматов без сжатия информации.}


   


   \KW{стеганография; стегоконтейнер; стегоанализ; детерминант; LSB; RS-сте\-го\-ана\-лиз; 


анализ с~использованием критерия хи-квад\-рат; анализ битовых слоев}





\DOI{10.14357\!/\!08696527200306}





\vskip 10pt plus 9pt minus 6pt





 \thispagestyle{headings}


 


 \vspace*{-12pt}


   


  \section{Введение}


  


  Стеганографическое встраивание~[1] стало перспективной областью 


информационной безопасности и~может использоваться для решения многих 


задач, например защиты авторского права, путем встраивания цифровой 


подписи, незаметной человеческому глазу. Помимо этого стеганографию 


можно использовать для передачи информации. 


  


  Одна из главных проблем стеганографии~--- поиск наиболее универсального 


и~при этом эффективного метода стеганографического встраивания, которое 


будет защищено от атак~[2].


  


  Максимальный объем данных, который можно встроить в~объект 


(стегоконтейнер) при условии, что его изменения не будут критическими, 


называется пропускной способностью в~контексте стеганографии.


  


  В настоящее время совершенно очевидно, что компьютерная стеганография 


стала одним из самых мощных инструментов скрытого встраивания 


защищаемых данных, которые могут храниться и~передаваться в~открытом 


незащищенном виде~\cite{3-lin, 4-lin, 5-lin}. Поэтому наибольший интерес для 


исследователей и~пользователей представляет совершенствование методов 


стеганографического встраивания скрываемой информации.


  


  На сегодняшний день при встраивании информации в~цветное изображение 


широко распространен метод на основе подмены наименее значащего бита 


(LSB)~[6, 7]. Основа данного метода заключается в~том, 


что при встраивании информации в~цветное изображение проводится замена от 


одного до четырех младших битов в~байтах пикселей изображения, что для 


человеческого глаза практически незаметно.


  


  У метода есть огромные преимущества в~объеме встраиваемой информации 


  и~быстродействии из-за его простоты~\cite{8-lin}. Однако простота метода не 


позволяет использовать его в~полной мере из-за его слабости к~отдельным 


методам стегоанализа


  


  \section{Постановка задачи}


  


  Если в~стеганографии в~качестве контейнера используется цветное 


изображение, то одним из вариантов отдельного элемента является пиксель. 


Рассмотрим цветовое пространство RGB. Каждый пиксель в~нем 


подразделяется на три цветовых области: красную (R~--- red), зеленую (G~--- green) и~синюю 


(B~--- blue), каждая из которых характеризуется отдельным числом от~0 до~255. 


  


  Эти числа можно представить в~двоичном коде с~использованием восьми 


знаков из нулей и~единиц. При замене последнего (младшего) бита разница при 


изменении цвета будет незаметна человеческому глазу. Благодаря этому 


обстоятельству можно добавлять информацию в~данные биты и~таким образом 


передавать ее~\cite{7-lin}.


  


  Метод при своей простоте и~эффективности имеет и~значительные минусы. 


На передачу информации при использовании алгоритма LSB могут повлиять 


самые незначительные искажения контейнера. При искажении контейнера 


меняется значение цвета, а~оно служит ключевой единицей при извлечении 


встроенной информации~\cite{9-lin}. По этой причине для данного метода 


подходят только изображения без сжатия информации. 


  


  Кроме того, он хорошо обнаруживается большинством стеганографических 


методов.


  


  Один из распространенных методов стегонализа~--- анализ битовых срезов. 


В~данном способе рассматриваются изображения, построенные с~по\-мощью 


значений отдельных битовых слоев~\cite{3-lin}.


  


  В ряде изображений четко различимы контуры объектов, поэтому 


дополнительная встроенная информация может быть заметна невооруженным 


глазом. На рис.~1 показан пример такого встраивания.


  


  


        \begin{figure} %fig1


        \vspace*{1pt}


 \begin{center}


 \mbox{%


 \epsfxsize=128mm 


\epsfbox{lin-1.eps}


 }


 \end{center}


\vspace*{-9pt}


  \Caption{Иллюстрация встроенной информации при использовании анализа битовых 


слоев: (\textit{а})~контейнер, в~который встроена информация методом LSB; 


(\textit{б})~битовый слой данного изображения (встроено 10\% информации от общего 


объема)}


  \end{figure}


  


  Как видно из данного примера, даже при анализе такими простыми методами 


информация уже становится заметна. В~более сложных случаях эффективные 


результаты дают статистические методы стегоанализа, такие как стегоанализ на 


основе критерия хи-квад\-рат и~RS-сте\-го\-ана\-лиз~\cite{4-lin, 5-lin}.


  %


  Поэтому при разработке нового способа необходимо учесть приведенные 


особенности и~недостатки, а~также обратить внимание на контейнер, который 


при использовании цветных изображений может иметь существенно разные 


исходные данные, влияющие на результат сокрытия информации.


  


  Отдельные попытки усовершенствования метода LSB делались  


в~\cite{11-lin, 12-lin}. Выбранное направление оказалось интересным, однако 


полученные результаты не давали полной уверенности в~надежности сокрытия 


информации. Поэтому у~исследователей остался вопрос относительно 


возможности усовершенствования метода по части повышения степени 


защищенности.


  


  \section{Алгоритм разработанного способа  }


  


  Следует напомнить, что разработанный способ берет за основу  


LSB-ал\-го\-ритм. Суть данного предложения заключается в~следующем.


  


  Выбранный контейнер цветного изображения делится на квадратные блоки 


одинакового размера. Для более детальной обработки лучше выбирать блоки 


минимального размера, например $2\times2$ или $3\times3$~пикселя. 


Определяется значение пикселя в~каждом блоке. При этом удобно использовать 


известный подход, описанный в~\cite{10-lin}, с~принятыми обозначениями. 


Можно задаться и~собственным представлением с~возможностью применения 


ко всем пикселям каждого блока. Далее в~каждом блоке вычисляется значение 


модуля детерминанта, которое связано с~координатами этого блока 


в~используемом контейнере. Затем выполняется ранжирование получившихся 


значений модуля детерминанта с~сохранением координат их расположения.


  


  Встраивание, в~свою очередь, выполняется в~те блоки, в~которых значение 


модуля детерминанта максимально. Значение модуля детерминанта 


используется для того, чтобы встраивание в~самые яркие или самые темные 


области проводилось в~последнюю очередь либо вообще не проводилось. 


Использование такого подхода маскирует контейнер. Этим ограничивается 


объем встраиваемой информации.


  


  Пошаговый алгоритм можно представить следующим образом.


  


  \smallskip


  


  \textit{Встраивание}:


  


  \smallskip


  


 \textbf{Шаг~1.} Разбиение пикселей изображения контейнера на блоки~$A$ 


определенного размера (например, $2\times2$). 


  


  Число зон для случая разбиения на блоки $2\times2$ пикселей вычисляется 


по следующей формуле:


  \begin{equation}


  n=\left\lfloor \fr{w}{2}\right\rfloor \cdot \left\lfloor \fr{h}{2}\right\rfloor\,,


  \end{equation}


гдe $w$ и~$h$~--- ширина и~высота изображения в~пикселях; $\lfloor\,\rfloor$~--- 


округление до ближайшего целого.


  


  \textbf{Шаг~2.} Значение пикселя в~зоне вычисляется, например, по формуле, 


приведенной в~\cite{10-lin}:


  \begin{equation}


  a_{i,j}=0{,}299 {R}+0{,}587 {G}+0{,}114 {B}\,,


  \end{equation}


где $i\in 0\ldots1$, $j\in 0\ldots 1$~--- координаты пикселя в~блоке;


${R}\hm\in 0\ldots 255$, ${G}\hm\in 0\ldots 255$ и~${B}\hm\in 


0\ldots 255$~--- цветовые компоненты пикселя.


  


  \textbf{Шаг~3.} Для каждой зоны вычисляется значение детерминанта


  \begin{equation*}


  \mathrm{det}\left\vert  D_k\right\vert\,,


  \end{equation*}


где $k\in 0\ldots n$.





  \textbf{Шаг~4.} Далее проводится ранжирование значений $D_k$ по убыванию 


с~сохранением координат.


  


  \textbf{Шаг~5.} Встраивание конфиденциальной информации происходит 


в~отсортированные блоки~$A$, которые имеют максимальное значение~$D_k$ 


в~три компоненты цвета: в~красную, синюю и~зеленую, т.\,е.\ на каждый 


пиксель приходится~3~бита информации. В~блок размера $2\times2$~пикселя 


встраивается~12~бит информации. 


  


  \smallskip


  


  


  Таким образом, максимальная информация, которую можно встроить 


в~изоб\-ра\-же\-ние для блока размером $2\times2$ пикселя, вычисляется 


следующим образом: 


  \begin{equation*}


  V=\left( \left\lfloor \fr{w}{2}\right\rfloor \cdot \left\lfloor \fr{h}{2}\right\rfloor 


\cdot 3\right) \cdot 2^2\,.


  \end{equation*}


  


  Исходя из полученного объема возможного встраивания информации, можно 


скорректировать его или выбрать необходимый, чтобы получить бо$\acute{\mbox{о}}$льшую 


защищенность при стегоанализе. 


  


  \smallskip


  


  


  \textit{Извлечение} происходит в~обратном порядке.


  


  \smallskip


  


  


  \textbf{Шаг~1.} Разбиение пикселей изображения на блоки~$A$ размером $2\times2$. 


Число блоков вычисляется по формуле~(1).


  


 \textbf{Шаг~2.} Значение пикселя в~зоне вычисляется по формуле~(2).


  


 \textbf{Шаг~3.} Для каждой зоны вычисляется значение дискриминанта по 


выбранной формуле.


  


 \textbf{Шаг~4.} Сортировка значений~$D_k$ по убыванию. 


  


 \textbf{Шаг~5.} Последовательно из блока~$D_k$ извлекается встроенное число 


битов информации.


  


  \section{Численное моделирование}


  


    \begin{floatingfigure}{60mm}


  {\small


  \begin{center}


  \begin{tabular}{|c|c|}


  \multicolumn{2}{c}{Ранжирование группы изображений}\\


  \multicolumn{2}{c}{\ }\\[-1pt]


  \hline


\tabcolsep=0pt\begin{tabular}{c}№\\ изображения\end{tabular}&\tabcolsep=0pt\begin{tabular}{c}Сумма площадей\\ пикселей области/\\


общую сумму\\ площадей пикселей 


\end{tabular}\\


\hline


1&0\hphantom{,}\\


2&\hphantom{,}$0{,}5\cdot 10^{-5}$\\


3&\hphantom{9}$33\cdot 10^{-5}$\\


4&$318\cdot 10^{-5}$\\


5&$334\cdot 10^{-5}$\\


6&$334\cdot 10^{-5}$\\


7&$557\cdot 10^{-5}$\\


8&$810\cdot 10^{-5}$\\


9&$989\cdot 10^{-5}$\\


10\hphantom{9}&$999\cdot 10^{-5}$\\


\hline


\end{tabular}


\end{center}


\vspace*{14pt}


}


\end{floatingfigure}


  


  Для проведения численного моделирования был выбран идентификационный 


параметр, указывающий отличие одного цветного изображения от другого. На 


сегодняшний день существует множество подобных  


параметров~\cite{1-lin, 11-lin}, однако каждый исследователь вправе выбрать 


такой, который может удовлетворять его требованиям. По этому параметру 


в~таб\-ли\-це приведено ранжирование выбранных цветных изображений. 


  


 





  


  Выбрана группа из десяти различных цветных изображений, каждое 


размером $1024\times1280$~пикселей (рис.~2). Встраивание различного объема 


информации проводилось в~каждое цветное изображение. Затем оценивались 


параметры различных способов стегоанализа.


  


 


  


  Рассмотрим использование стегоатаки на основе анализа битовых срезов. На 


рис.~3 показаны битовые срезы контейнеров оригинального LSB-ме\-то\-да 


и~разработанного способа при условии, что в~них встроен одинаковый объем 


информации. 


    


    


  


  На рис.~3 видно, что разработанная модификация позволяет скрыть 


информацию при использовании данного метода стегоанализа. 


  


  Рассмотрим результаты других широко используемых методов стегоанализа.


  


  RS-стегоанализ относится к~статистическим методам стегоанализа. 


Изображение разбивается на группы пикселей. Деление на группы происходит 


после применения флип\-пинг-про\-це\-ду\-ры с~применением введенной маски  


и~дис\-кри\-ми\-нант-функ\-ции. После проведения расчетов анализируются 


значения дис\-кри\-ми\-нант-функ\-ции до и~после применения процедуры 


флиппинга. 





 \setcounter{figure}{1}


  \begin{figure} %fig2


   \vspace*{1pt}


 \begin{center}


 \mbox{%


 \epsfxsize=128mm 


\epsfbox{lin-2.eps}


 }


 \end{center}


\vspace*{-9pt}


  \Caption{Группа изображений}


  \end{figure}


  


 


  


  Работа этого алгоритма строится на предположении равенства регулярных 


и~сингулярных групп в~пустом изображении после применения процедуры 


флиппинга. Иначе считается, что в~изображение встроена информация.





 \begin{figure} %fig3


   \vspace*{1pt}


 \begin{center}


 \mbox{%


 \epsfxsize=128mm 


\epsfbox{lin-3.eps}


 }


 \end{center}


\vspace*{-9pt}


  \Caption{Битовый срез контейнеров с~10\% внедренной информации: 


(\textit{a})~стандартный метод LSB; (\textit{б})~модифицированный метод LSB}


  \end{figure}


  


  На рис.~4 показаны графики значений RS-сте\-го\-ана\-ли\-за 


для~10~изображений. Значение RS в~данном способе коррелирует с~процентом 


измененной информации~[2].





\begin{figure} %fig4


 \vspace*{1pt}


 \begin{center}


 \mbox{%


 \epsfxsize=128mm 


\epsfbox{lin-4.eps}


 }


 \end{center}


 \vspace*{-9pt}


\begin{minipage}[t]{64mm}


\Caption{Результаты RS-стегоанализа для стандартного LSB~(\textit{1}) и~для разработанной 


модификации~(\textit{2}): (\textit{а})~при заполнении 10\% объема контейнера; (\textit{б})~при 


заполнении 30\% объема контейнера; (\textit{в})~при заполнении 50\% объема контейнера}


\end{minipage}


%\end{figure}


\hfill


% \begin{figure} %fig5


\begin{minipage}[t]{64mm}


  \Caption{Результаты стегоанализа с~использованием критерия хи-квадрат для 


стандартного LSB~(\textit{1}) и~разработанной модификации~(\textit{2}): (\textit{а})~при заполнении 10\% объема 


контейнера; (\textit{б})~при заполнении 30\% объема контейнера; (\textit{в})~при 


заполнении 50\% объема контейнера}


\end{minipage}


\end{figure}


  


 


  


  Как видно из графиков, встраивание малого объема информации имеет 


практически нулевые значения RS, что, в~свою очередь, может 


свидетельствовать о защите от использования данного метода, так как 


атакующим полученные результаты могут быть восприняты как отсутствие 


встроенной информации или как шум.


  


  Следующий способ стегоанализа основан на применении критерия  


хи-квад\-рат~[5]. Он также относится к~статистическим видам стегоанализа 


и~оценивает предположение о том, что в~пустом контейнере распределение цвета 


происходит таким образом, что вероятность одновременного появления 


соседних цветов, т.\,е.\ тех, которые различаются на один младший бит, 


незначительна. В~противном случае делается вывод о том, что в~контейнере 


находятся посторонние данные. 


  


  На рис.~5 представлены графики, показывающие значения распределения  


хи-квад\-рат на изображениях, полученных при использовании стандартного 


встраивания методом LSB и~с применением разработанного алгоритма, для 


случаев встраивания разного объема информации.


  


 


  


  Из графиков, приведенных на рис.~5, видно, что значение вероятности для 


предложенного способа равно нулю на всем промежутке. Таким образом, 


можно сделать вывод о том, что способ является защищенным к~применению 


атак на основе критерия хи-квад\-рат.


  


  \section{Обсуждение результатов и~выводы}


  


  По приведенным результатам численного моделирования можно сделать 


предварительный вывод о том, что предложенный статический способ 


сокрытия информации достаточно хорошо справляется с~защитой от 


стеганографических атак с~использованием широко известных способов 


стегоанализа. 


  


  Эксперименты показали, что предложенный способ полностью защищен от 


метода анализа битовых срезов и~от анализа с~использованием критерия  


хи-квад\-рат, RS-сте\-го\-ана\-ли\-за, анализа отношения цветов при условии, 


что происходит встраивание небольших объемов информации.


  


  В отличие от других модификаций метода LSB предложенный способ более 


универсален и~позволяет учесть особенности конкретного контейнера.


  %


  Кроме того, предложенный способ позволяет встраивать б$\acute{\mbox{о}}$льшие объемы 


информации, чем классический LSB.


  


  Необходимо отметить, что эффективность работы алгоритма зависит от 


особенностей конкретного цветного изображения. Если в~изображении есть 


большие области слишком ярких либо слишком темных цветов, то существует 


опасность возникновения такой ситуации, при которой информация будет 


записываться подряд, что понизит стеганостойкость.


  


{\small\frenchspacing


{%\baselineskip=10.8pt


\addcontentsline{toc}{section}{Литература}


\begin{thebibliography}{99}


\bibitem{1-lin}


\Au{Коханович Г.\,Ф., Пузыренко~А.\,Ю.} Компьютерная стеганография. Теория 


и~практика.~--- Киев: МК-Пресс, 2006. 288~с.


\bibitem{2-lin}


\Au{Абазина Е.\,С., Ерунов А.\,А.} Цифровая стеганография: состояние и~перспективы~/\!\!/ Сис\-те\-мы 


управ\-ле\-ния, связи и~безопас\-ности, 2016. №\,2. С.~182--201.





\bibitem{5-lin} %3


\Au{Westfield A., Pfitzmann~A.} Attacks on steganographic systems: Breaking the steganographic 


utilities EzStego, Jsteg, Steganos, and S-Tools~--- and some lessons learned~/\!\!/ 


  Information hiding~/ Ed. A.~Pfitzmann.~--- Lecture notes in 


computer Science ser.~--- Springer, 1999. Vol.~1768. P.~61--76.





\bibitem{4-lin} %4


\Au{Fridrich J., Goljan~M., Du~R.} Reliable detection of LSB steganography in color and 


grayscale images~/\!\!/ Workshop on Multimedia and Security: New Challenges 


Proceedings.~--- Binghamton, NY, USA: State University of New York, 2001. P.~27--30. doi: 


10.1145/1232454.1232466.





\bibitem{3-lin} %5


\Au{Грибунин В.\,Г., Оков~И.\,Н., Туринцев~И.\,В.} Цифровая стеганография.~---  


М.: Солон-Пресс, 2002. 272~с.








\bibitem{6-lin}


\Au{Fridrich J., Du~R., Meng~L.} Steganalysis of LSB encoding in color images~/\!\!/ IEEE  


Conference (International) on Multimedia and Expo.~--- IEEE Computer Society Press, 2000. 


Vol.~3. P.~1279--1282.


\bibitem{7-lin}


\Au{Шмаев В.\,Б.} Современная стеганография. Принципы, основные носители и~методы 


противодействия, 2005. {\sf  


https://docplayer.ru/29953095-Sovremennaya-steganografiya-principy-osnovnye-nositeli-i-metody-protivodeystviya.html.} 


\bibitem{8-lin}


\Au{Urbanovich N., Plaskovitsky~V.} The use of steganographic techniques for protection of 


intellectual property rights~/\!\!/ New Electrical and Electronic Technologies and Their Industrial 


Implementation: 7th  Conference (International).~--- Zakopane, Poland, 2011. P.~147--148.


\bibitem{9-lin} %9


\Au{Журавель И.\,М.} Краткий курс теории обработки изображений, 2019. {\sf 


http://matlab.exponenta.ru/imageprocess/book2/index.php}.





\bibitem{11-lin} %10


\Au{Кривошеев И.\,А., Линник~М.\,А.} К~вопросу об оценке устойчивости 


стеганографической системы~/\!\!/ Ученые заметки ТОГУ, 2017. Т.~8. №\,2. С.~433--437.


\bibitem{12-lin} %11


\Au{Кривошеев И.\,А., Линник~М.\,А.} К~вопросу об оценке пропускной способности 


стенографической системы~/\!\!/ Информационные технологии и~высокопроизводительные 


вычисления: Мат-лы IV Всеросс. научн.-практич. конф.~--- Хабаровск: ТОГУ, 2017.  


С.~94--97.





\bibitem{10-lin} %12


\Au{Гонсалес Р., Вудс~Р.} Цифровая обработка изображений.~--- М.: Техносфера, 2005. 


1072~с.





    \end{thebibliography}


} }








\vspace*{-6pt}





\hfill{\small\textit{Поступила в~редакцию 21.04.20}}














\vspace*{6pt}





\hrule





\vspace*{2pt}





%\newpage





%\vspace*{-28pt}





\hrule





\vspace*{8pt}





\def\tit{STATIC METHOD FOR~STEGANOGRAPHIC EMBEDDING\\ BASED ON~LEAST SIGNIFICANT BIT}





\def\titkol{Static method for~steganographic embedding based on~least significant bit}





\def\aut{I.\,A.~Krivosheev$^1$ and~M.\,A.~Linnik$^2$}





\def\autkol{I.\,A.~Krivosheev and~M.\,A.~Linnik}





\titel{\tit}{\aut}{\autkol}{\titkol}





\vspace*{-15pt}








 \noindent


$^1$Mining Institute of the Far Eastern Branch of the Russian Academy of Sciences,\linebreak


$\hphantom{^1}$Khabarovsk Federal 


Research Center of the Far Eastern Branch of the Russian\linebreak


$\hphantom{^1}$Academy of Sciences, 51~Turgenev Str., 


Khabarovsk 680000,Russian Federation





\noindent


$^2$Computing Center, Far Eastern Branch of the Russian Academy of Sciences,\linebreak


$\hphantom{^1}$65~Kim Yu Chen 


Str., Khabarovsk 680000, Russian Federation





\def\leftfootline{\small{\textbf{\thepage}


\hfill Sistemy i Sredstva Informatiki~---


Systems and Means of Informatics 2020 vol~30 no\,3}


}%


 \def\rightfootline{\small{Sistemy i Sredstva Informatiki~---


 Systems and Means of Informatics 2020 vol~30 no\,3


\hfill \textbf{\thepage}}}





\vspace*{3pt}











\Abste{A new approach to information hiding based on the LSB (least significant bit)


method is proposed. A static method for 


embedding information\linebreak\vspace*{-12pt}}





\Abstend{in a color image based on the matrix representation of individual blocks and 


estimating their determinant modulus is proposed. Numerical modeling evaluated the possibility of the 


proposed method to withstand various methods of steganalysis. Its advantage is shown both in the degree of 


security and in the amount of embedded information. This algorithm can be used for embedding in image by 


formats without compressing information.}








\KWE{steganography; stego container; steganalysis; determinant; LSB; RS-steganalysis; chi-squared test 


steganalysis; visual analysis of bit slices}





\DOI{10.14357\!/\!08696527200306}





%\vspace*{-24pt}





% \Ack


 % \noindent


 





\vspace*{-12pt}





\renewcommand{\bibname}{\protect\rmfamily References}


%\renewcommand{\bibname}{\large\protect\rm References}





{\small\frenchspacing


{%\baselineskip=10.8pt


\addcontentsline{toc}{section}{References}


\begin{thebibliography}{99}


\bibitem{1-lin-1}


\Aue{Kokhanovich, G.\,F., and A.\,Yu.~Puzyrenko.} 2006. \textit{Komp'yuternaya steganografiya. Teoriya 


i~praktika} [Computer steganography. Theory and practice]. Kiev: MK-Press. 288~p.


\bibitem{2-lin-1}


\Aue{Abazina, E.\,S., and A.\,A.~Erunov.} 2016. Tsifrovaya steganografiya: sostoyanie i~perspektivy [Digital steganography: 


Status and development outlook]. \textit{Sistemy uprav\-le\-niya, svyazi i~bezopasnosti} [Systems of


Control, Communication and Security] 2:182--201.





\bibitem{5-lin-1} %3


\Aue{Westfield, A., and A.~Pfitzmann.} 1999. Attacks on steganographic systems: Breaking the 


steganographic utilities EzStego, Jsteg, Steganos and S-Tools~--- and some lessons learned. \textit{Information 


hiding}. Ed.\ A.~Pfitzmann. Lecture notes in computer science  ser.  Springer. 1768:61--76. 








\bibitem{4-lin-1} %4


\Aue{Fridrich, J., M.~Goljan, and R.~Du.} 2001. Reliable detection of LSB steganography in color and 


grayscale images. \textit{Workshop on Multimedia and Security: New Challenges Proceedings}. 


Binghamton, NY: State Unoversity of New York. 27--30. doi: 10.1145\!/\!1232454.1232466.





\bibitem{3-lin-1} %5


\Aue{Gribunin, V.\,G., I.\,N.~Okov, and I.\,V.~Turintsev.} 2002. \textit{Tsifrovaya steganografiya} [Digital 


steganography]. Moscow: Solon-Press. 272~p.





\bibitem{6-lin-1}


\Aue{Fridrich, J., R.~Du, and L.~Meng.} 2000. Steganalysis of LSB encoding in color images. 


\textit{Conference (International) on Multimedia and Expo}. IEEE Computer Society Press. 3:1279--1282.


\bibitem{7-lin-1}


\Aue{Shmaev, V.\,B.} 2005. Sovremennaya steganografiya. Printsipy, osnovnye nositeli i~metody 


protivodeystviya [Modern steganography. Principles, main carriers and methods of counteraction]. Available 


at: {\sf 


 https://docplayer.ru/29953095-Sovremennaya-steganografiya-principy-osnovnye-nositeli-i-metody-protivodeystviya.html} 


(accessed April~22, 2020).


\bibitem{8-lin-1}


\Aue{Urbanovich, N., and V.~Plaskovitsky.} 2011. The use of steganographic techniques for protection of 


intellectual property rights. \textit{New Electrical and Electronic Technologies and Their Industrial 


Implementation: 7th  Conference (International)}. Zakopane, Poland. 147--148.


\bibitem{9-lin-1}


\Aue{Zhuravel', I.\,M.} 2019. Kratkiy kurs teorii obrabotki izobrazheniy [Image processing toolbox]. 


Available at: {\sf  


http://matlab.exponenta.ru/imageprocess/book2/index.php} (accessed April~22, 2020).





\bibitem{11-lin-1} %10


\Aue{Krivosheev, I.\,A., and M.\,A.~Linnik.} 2017. K voprosu ob otsenke ustoychivosti 


steganograficheskoy sistemy [On the issue of assessing the stability of a steganographic system]. 


\textit{Uchenye zametki TOGU} [TOGU Science Notes] 8(2):433--437. 


\bibitem{12-lin-1} %11


\Aue{Krivosheev, I.\,A., and M.\,A.~Linnik.} 2017. K~voprosu ob otsenke propusknoy sposobnosti 


stenograficheskoy sistemy [On the issue of assessing the bandwidth of a shorthand system]. 


\textit{Informatsionnyye tekhnologii i~vysokoproizvoditel'nye vychisleniya: mat-ly IV Vseross.  


nauchn.-praktich. konf.} [Information Technology and High Performance Computing: Materials of the 4th 


All-Russian Scientific and Practical Conference]. Khabarosk: TOGU. 94--97.


\bibitem{10-lin-1} %12


\Aue{Gonsales, R., and R.~Vuds.} 2005. \textit{Tsifrovaya obrabotka izobrazheniy} [Digital image 


processing]. Moscow: Tekhnosfera. 1072~p.


\end{thebibliography}


} }





%\vspace*{-3pt}





\hfill{\small\textit{Received April 21, 2020}} 





%\vspace*{-16pt} 





\Contr





\noindent


\textbf{Krivosheev Igor A.} (b.\ 1949)~--- Doctor of Science in technology, leading scientist, Mining 


Institute of the Far Eastern Branch of the Russian Academy of Sciences, Khabarovsk Federal Research 


Center of the Far Eastern Branch of the Russian Academy of Sciences, 51~Turgenev Str., Khabarovsk 


680000,Russian Federation; \mbox{igork@as.khb.ru}





\vspace*{3pt}





\noindent


\textbf{Linnik Maxim A.} (b.\ 1995)~--- junior scientist, Computing Center, Far Eastern Branch of the 


Russian Academy of Sciences, 65~Kim Yu Chen Str..Khabarovsk 680000, Russian Federation; 


\mbox{linnik.max1995@mail.ru}





\label{end\stat}





\renewcommand{\bibname}{\protect\rm Литература} 


